% Generated draft: requires mathematical and source review.
\documentclass[11pt]{article}
\usepackage[margin=25mm]{geometry}
\usepackage{fontspec}
\setmainfont{Latin Modern Roman}
\usepackage{amsmath,amssymb,mathtools,unicode-math}
\setmathfont{Latin Modern Math}
\usepackage{enumitem,needspace,xurl,hyperref}
\setlength{\parindent}{0pt}
\setlength{\parskip}{5pt}
\setlength{\emergencystretch}{3em}
\newcommand{\N}{\mathbb{N}}
\newcommand{\Z}{\mathbb{Z}}
\newcommand{\Q}{\mathbb{Q}}
\newcommand{\R}{\mathbb{R}}
\newcommand{\C}{\mathbb{C}}
\newcommand{\F}{\mathbb{F}}
\newcommand{\A}{\mathbb{A}}
\newcommand{\PP}{\mathbb{P}}
\newcommand{\E}{\mathbb{E}}
\newcommand{\eps}{\varepsilon}
\newcommand{\dd}{\,\mathrm{d}}
\begin{document}
\section*{UnsolvedMath ID: 30002093}
\textbf{Source problem code:} catalogue entry
\subsection{Definitions and mathematical statement}
Let \(X\) be a Banach space and let \(M\) be a positive integer. Write
\[
\mathcal S_M(X)=\{(V_1,\ldots,V_M): V_j\subseteq X\text{ is a closed linear subspace for each }j\}.
\]
The source asks for necessary and sufficient conditions for the subspace-segmentation approximation error with exponent \(p\ne2\) to attain its minimum over such tuples.

If \(\mathcal E_p\) denotes the intended approximation-error functional, the requested attainment property has the schematic form
\[
\exists\,\mathbf V^*\in\mathcal S_M(X)\quad
\forall\,\mathbf V\in\mathcal S_M(X),\qquad
\mathcal E_p(\mathbf V^*)\leq\mathcal E_p(\mathbf V).
\]
Here \(\mathcal E_p\) is an unresolved placeholder, not a functional defined by the supplied record. The record specifies neither the data being approximated nor the formula for their approximation error, the allowed range of \(p\) beyond \(p\ne2\), or any restrictions on the dimensions or admissibility of the subspaces. In particular, the displayed family includes \((X,\ldots,X)\); whether the intended segmentation problem excludes this tuple is unspecified. If a restricted admissible family is intended, that family must replace \(\mathcal S_M(X)\) in the attainment condition.

Thus the supplied question is to characterize minimizer existence for the intended problem when \(p\ne2\), but a fully specified mathematical statement requires the missing objective, input data, admissibility hypotheses, and quantifier scope. The record does not determine whether conditions are sought for a particular data set or uniformly over a class of data sets.
\subsection{Short English statement}
The problem asks exactly when an optimal collection of \(M\) closed subspaces exists for Banach-space subspace segmentation with exponent \(p\ne2\). The supplied record leaves the approximation error, data, exponent range, and subspace restrictions unspecified, so the displayed attainment condition is only a statement schema requiring human mathematical review.
\subsection{Sources}
\begin{itemize}
\item Local UnsolvedMath registry: permanent ID \(30002093\), title ``Minimizers for Banach-Space Subspace Segmentation,'' source problem code \texttt{catalogue entry}, collection ``Oberwolfach Reports Open Problems.'' No individual proposer is supplied; the recorded proposed year is 2012.
\item Supplied report citation: \emph{Applied Harmonic Analysis and Sparse Approximation (2013)}, Oberwolfach Reports, DOI 10.4171/owr/2012/29, pp.~1759--1843. \url{https://doi.org/10.4171/owr/2012/29}. The background describes these reports as documenting mathematical workshops and their open-problem discussions. The citation's displayed year is 2013, whereas the DOI and recorded proposed year refer to 2012; their relationship is not explained in the record.
\item Supplied literature reference: Aldroubi et al., ``review of subspace segmentation.'' \url{https://onlinelibrary.wiley.com/doi/10.1155/2013/417492}. The record supplies no fuller bibliographic details or precise supporting passage.
\item Supplied literature assessment, checked on \texttt{2026-08-21}: the record labels the problem open and states that necessary and sufficient existence conditions are known for \(p=2\), while the cited survey and subsequent literature do not provide a general characterization for \(p\ne2\). This is supplied dated curation derived from the OWR record, not an independently verified literature assessment. The supplied background explicitly calls for independent verification of time-sensitive status and cited claims.
\end{itemize}
\end{document}
